\documentclass[11pt,a4paper]{article}

% --- Encoding ---
\usepackage[utf8]{inputenc}
\usepackage[T1]{fontenc}
\usepackage[spanish,es-noshorthands]{babel}
\usepackage{lmodern}

% --- Layout ---
\usepackage[margin=2.5cm]{geometry}
\usepackage{parskip}
\setlength{\parindent}{0pt}
\setlength{\parskip}{0.5em plus 0.2em minus 0.1em}

% --- Tables and graphics ---
\usepackage{booktabs}
\usepackage{array}
\usepackage{enumitem}
\usepackage{xcolor}
\usepackage{graphicx}
\usepackage{tikz}
\usetikzlibrary{positioning,shapes.geometric,arrows.meta,fit,backgrounds,calc,decorations.pathreplacing}
\usepackage{caption}
\usepackage{subcaption}
\usepackage{amsmath,amssymb}

% --- Soporte para caracteres Unicode usados en formulas en texto plano ---
\usepackage{newunicodechar}
\newunicodechar{→}{\ensuremath{\rightarrow}}
\newunicodechar{≥}{\ensuremath{\geq}}
\newunicodechar{≤}{\ensuremath{\leq}}
\newunicodechar{×}{\ensuremath{\times}}
\newunicodechar{÷}{\ensuremath{\div}}
\newunicodechar{≈}{\ensuremath{\approx}}
\newunicodechar{⌈}{\ensuremath{\lceil}}
\newunicodechar{⌉}{\ensuremath{\rceil}}
\newunicodechar{✓}{\ensuremath{\checkmark}}

% --- Code listings ---
\usepackage{listings}
\definecolor{codegray}{rgb}{0.96,0.96,0.96}
\definecolor{codeborder}{rgb}{0.85,0.85,0.85}
\definecolor{codecomment}{rgb}{0.30,0.55,0.30}
\definecolor{codekeyword}{rgb}{0.10,0.30,0.70}
\definecolor{codestring}{rgb}{0.65,0.15,0.10}

\lstdefinestyle{cstyle}{
    language=C,
    backgroundcolor=\color{codegray},
    basicstyle=\ttfamily\footnotesize,
    breaklines=true,
    frame=single,
    rulecolor=\color{codeborder},
    showstringspaces=false,
    keywordstyle=\color{codekeyword}\bfseries,
    commentstyle=\color{codecomment}\itshape,
    stringstyle=\color{codestring},
    aboveskip=1em,
    belowskip=1em,
}

\lstdefinestyle{shellstyle}{
    language=bash,
    backgroundcolor=\color{codegray},
    basicstyle=\ttfamily\small,
    breaklines=true,
    frame=single,
    rulecolor=\color{codeborder},
    showstringspaces=false,
    keywordstyle=\color{codekeyword}\bfseries,
    commentstyle=\color{codecomment}\itshape,
    stringstyle=\color{codestring},
    aboveskip=1em,
    belowskip=1em,
    columns=fullflexible,
    upquote=true,
}

\lstdefinestyle{textstyle}{
    backgroundcolor=\color{codegray},
    basicstyle=\ttfamily\footnotesize,
    breaklines=true,
    frame=single,
    rulecolor=\color{codeborder},
    showstringspaces=false,
    aboveskip=1em,
    belowskip=1em,
    columns=fullflexible,
}

\lstset{style=textstyle}

% --- Definition / Idea / Important boxes ---
\usepackage[most]{tcolorbox}
\newtcolorbox{idea}[1][]{
    colback=orange!5!white,
    colframe=orange!60!black,
    fonttitle=\bfseries,
    title=Idea intuitiva,
    sharp corners=southwest,
    rounded corners=northwest,
    arc=2pt,
    boxrule=0.6pt,
    left=8pt,right=8pt,top=4pt,bottom=4pt,
    #1
}
\newtcolorbox{importante}[1][]{
    colback=red!5!white,
    colframe=red!50!black,
    fonttitle=\bfseries,
    title=Importante,
    sharp corners=southwest,
    rounded corners=northwest,
    arc=2pt,
    boxrule=0.6pt,
    left=8pt,right=8pt,top=4pt,bottom=4pt,
    #1
}
\newtcolorbox{respuesta}[1][]{
    colback=green!5!white,
    colframe=green!40!black,
    fonttitle=\bfseries,
    title=Respuesta,
    sharp corners=southwest,
    rounded corners=northwest,
    arc=2pt,
    boxrule=0.6pt,
    left=8pt,right=8pt,top=4pt,bottom=4pt,
    #1
}

% --- Hyperlinks ---
\usepackage[colorlinks=true,
            linkcolor=blue!60!black,
            urlcolor=blue!50!black,
            citecolor=blue!60!black]{hyperref}

% --- Color palette ---
\definecolor{c1}{HTML}{4C72B0}
\definecolor{c2}{HTML}{DD8452}
\definecolor{c3}{HTML}{55A868}
\definecolor{c4}{HTML}{8172B2}
\definecolor{cred}{HTML}{C44E52}
\definecolor{cfree}{HTML}{8C8C8C}

% --- Title block ---
\title{
    \vspace{-2.5em}
    \textbf{Parcial --- OS} \\[0.3em]
    \large Sistemas Operativos
}
\author{
    Tomás Rodeghiero \\
    \small UNRC --- Universidad Nacional de Río Cuarto
}
\date{}

\begin{document}

\maketitle

\begin{abstract}
\noindent
Resolución completa y verificada del parcial ``Parcial --- OS'' de
la cátedra de Sistemas Operativos. El parcial consta de diez
ejercicios mixtos (multiple-choice y abiertos) sobre: diseño de
sistemas operativos (monolítico vs.\ microkernel), gestión de
procesos, motivación de \texttt{fork+exec}, características de xv6,
deadlock en pares de locks adquiridos en orden inverso, traducción
de direcciones lógicas a físicas con paginación de dos niveles,
diagrama de tabla de páginas para un proceso xv6, segmentación vs.\
paginación, paginación de un solo nivel en x86, y concepto de
memoria virtual. Cada respuesta se verifica de manera independiente
y se documenta el razonamiento.
\end{abstract}

\tableofcontents
\bigskip

% =====================================================================
\section{Ejercicio 1 --- Diseño de sistemas operativos}
% =====================================================================

\textbf{Enunciado.} Marcar las respuestas verdaderas.

\subsection*{1.a --- ``Un sistema monolítico tiene los subsistemas del kernel en un mismo espacio de memoria''}

\begin{respuesta}
\textbf{V}. En un kernel monolítico (Linux, Windows NT, xv6) todos
los subsistemas (scheduler, gestor de memoria, FS, drivers,
networking) se compilan en un único binario y se cargan en el mismo
espacio de direcciones del kernel. La comunicación entre subsistemas
es por llamadas a función directa y acceso a estructuras globales.
\end{respuesta}

\subsection*{1.b --- ``En un sistema monolítico sus subsistemas se comunican por medio de mensajes''}

\begin{respuesta}
\textbf{F}. En monolítico, los subsistemas comparten el mismo
espacio de direcciones y se comunican por \textbf{llamadas a función
directa} y por acceso a estructuras de datos globales (tabla de
procesos, tabla de archivos, etc.). El paso de mensajes es típico
del \emph{microkernel}, no del monolítico.
\end{respuesta}

\subsection*{1.c --- ``En un microkernel un subsistema no puede acceder al espacio de memoria de otros módulos''}

\begin{respuesta}
\textbf{V}. En un microkernel los subsistemas (FS, drivers, etc.)
son \emph{procesos} en modo usuario, cada uno con su propio espacio
de direcciones. La MMU impide que un servicio acceda directamente a
la memoria de otro. La comunicación es vía IPC (mensajes) gestionado
por el microkernel.
\end{respuesta}

\subsection*{1.d --- ``Un diseño microkernel tiene como ventaja que se puede distribuir\ldots''}

\emph{(diferentes componentes corriendo en diferentes computadoras
en una red)}

\begin{respuesta}
\textbf{V}. Como los servicios del SO son procesos independientes
que se comunican por mensajes, esos mensajes pueden viajar tanto
localmente (IPC tradicional) como a través de la red. Esto fue una
ventaja explotada por Mach (base de macOS) y por sistemas
distribuidos como QNX/Helios. Un servicio (e.g., un servidor de
archivos) puede vivir físicamente en otra máquina y los clientes lo
acceden de manera transparente.

\textbf{En un monolítico esto es imposible}: todos los subsistemas
viven en el mismo proceso (el kernel) y dependen de acceso directo
a la memoria, lo cual no es distribuible.
\end{respuesta}

\subsection*{1.e --- ``Un sistema microkernel generalmente ofrece mejor rendimiento que un monolítico''}

\begin{respuesta}
\textbf{F}. Es al revés: \emph{en general}, los monolíticos tienen
\textbf{mejor rendimiento}. La razón: en monolítico una llamada
entre subsistemas es una llamada a función directa; en microkernel
es un mensaje IPC con (al menos) dos cambios de modo y dos cambios
de contexto, más posible copia de datos entre espacios.

A cambio, el microkernel gana en \textbf{modularidad}, \textbf{robustez}
(crash de un driver no tira el kernel), \textbf{seguridad} y
\textbf{distribución} (1.d).
\end{respuesta}

\subsection*{Resumen Ej.\ 1}

\begin{center}
\small
\begin{tabular}{@{}lll@{}}
\toprule
Ítem & Respuesta & Marcar como verdadera \\
\midrule
a    & V         & \checkmark \\
b    & F         &           \\
c    & V         & \checkmark \\
d    & V         & \checkmark \\
e    & F         &           \\
\bottomrule
\end{tabular}
\end{center}

\begin{importante}
\textbf{Corrección frente a la respuesta original}: el ítem
\textbf{(d) es V}. La capacidad de distribuir un microkernel en una
red es justamente una de sus ventajas clásicas.
\end{importante}

% =====================================================================
\section{Ejercicio 2 --- Gestión de procesos}
% =====================================================================

\textbf{Enunciado.} Marcar las respuestas verdaderas.

\subsection*{2.a --- ``En UNIX la única llamada al sistema que crea un nuevo proceso es \texttt{fork()}''}

\begin{respuesta}
\textbf{V}. \texttt{fork()} es la única syscall que crea un nuevo
proceso (con un nuevo PID). Sus variantes \texttt{vfork()} y
\texttt{clone()} en Linux son del mismo mecanismo conceptual:
producen un nuevo proceso. \texttt{exec()}, en contraste, NO crea
proceso: reemplaza la imagen del proceso actual.
\end{respuesta}

\subsection*{2.b --- ``Un proceso \emph{zombie} es uno tal que no puede finalizar porque no ha liberado todos sus recursos''}

\begin{respuesta}
\textbf{F}. Un zombie \emph{ya finalizó} (ejecutó \texttt{exit()} o
fue terminado por una señal). Sus recursos (memoria, descriptores,
mapeos) \textbf{ya fueron liberados}. Lo único que sigue existiendo
es su \emph{descriptor} (PCB) en la tabla de procesos del kernel,
con su \emph{exit code} guardado, esperando que el \textbf{padre}
llame a \texttt{wait()}/\texttt{waitpid()} para recogerlo.

Sólo cuando el padre hace \texttt{wait}, el descriptor se libera y
el zombie desaparece definitivamente.
\end{respuesta}

\subsection*{2.c --- ``La llamada al sistema \texttt{exec(path, argv)} crea un nuevo proceso y carga el código y datos estáticos en memoria''}

\begin{respuesta}
\textbf{F}. \texttt{exec()} \textbf{no crea un nuevo proceso}.
Reemplaza la imagen del proceso \emph{corriente}: mismo PID, mismo
descriptor, mismo padre, pero todo el código, datos y stack se
sustituyen por los del nuevo programa.

Para crear procesos se usa \texttt{fork()}. El patrón clásico es
\texttt{fork} + \texttt{exec}: \texttt{fork} crea el hijo;
\texttt{exec} reemplaza al hijo con el nuevo programa.
\end{respuesta}

\subsection*{2.d --- ``Siempre hay un proceso en estado RUNNING, ya que el scheduler asume eso''}

\begin{respuesta}
\textbf{F}. No es cierto que \emph{siempre} haya un proceso de
usuario en RUNNING:

\begin{itemize}[leftmargin=*,nosep]
    \item En xv6 (original), si la cola \texttt{RUNNABLE} está vacía,
        el \texttt{scheduler()} entra a un \emph{spin} (loop) con la
        instrucción \texttt{hlt} o equivalente, esperando que un
        \texttt{wakeup} o el timer cree trabajo. No hay un ``proceso
        idle'' formal.
    \item En Linux y algunos SO modernos, sí existe un \emph{idle
        process} (\texttt{swapper} con PID 0) que se planifica
        cuando no hay nadie más. Pero eso es una particularidad de
        la implementación, no un requisito conceptual.
\end{itemize}

El scheduler está \emph{diseñado para tolerar} el caso de no haber
procesos RUNNABLE. No \emph{asume} que siempre haya algo que correr.
\end{respuesta}

\subsection*{2.e --- ``Sólo puede haber un único proceso en estado RUNNING''}

\begin{respuesta}
\textbf{F}. En sistemas multi-procesador (SMP), hay tantos procesos
RUNNING simultáneamente como CPUs/cores haya. Cada CPU ejecuta a su
proceso elegido por el scheduler. En sistemas mono-procesador sí
sería V, pero la afirmación absoluta es falsa.
\end{respuesta}

\subsection*{Resumen Ej.\ 2}

\begin{center}
\small
\begin{tabular}{@{}lll@{}}
\toprule
Ítem & Respuesta & Marcar como verdadera \\
\midrule
a    & V         & \checkmark \\
b    & F         &           \\
c    & F         &           \\
d    & F         &           \\
e    & F         &           \\
\bottomrule
\end{tabular}
\end{center}

\begin{importante}
\textbf{Sólo la (a) debe marcarse como verdadera.} La respuesta
original tenía sin completar (d) y (e); ambas son \textbf{F}.
\end{importante}

% =====================================================================
\section{Ejercicio 3 --- Ventajas de \texttt{fork()} + \texttt{exec()}}
% =====================================================================

\textbf{Enunciado.} Describa las ventajas de que ejecutar un programa
en el contexto de un nuevo proceso deba hacerse en dos pasos
(\texttt{fork()} \ldots\ \texttt{exec()}). Ejemplificar.

\subsection*{Análisis}

La división en dos syscalls (\texttt{fork}, \texttt{exec}) es una
\textbf{decisión de diseño elegante} de UNIX que evita un único
\texttt{spawn(prog, args, \ldots)} con docenas de argumentos. La
idea: entre \texttt{fork} y \texttt{exec} el hijo está en una
\emph{ventana} en la que puede usar \emph{cualquier syscall}
disponible para configurar su ambiente, y luego saltar al programa
final.

\subsection*{Ventajas concretas}

\paragraph{1) Redirecciones de E/S por el shell.}
El shell hace \texttt{open}/\texttt{close}/\texttt{dup2} entre
\texttt{fork} y \texttt{exec} para conectar \texttt{stdin}/\texttt{stdout}/
\texttt{stderr} a archivos o pipes:

\begin{lstlisting}[style=cstyle]
/* Implementacion de `ls > out.txt` por el shell */
pid_t pid = fork();
if (pid == 0) {
    int fd = open("out.txt", O_WRONLY | O_CREAT | O_TRUNC, 0644);
    dup2(fd, STDOUT_FILENO);   /* fd 1 ahora apunta a out.txt */
    close(fd);
    execlp("ls", "ls", (char*)NULL);
}
\end{lstlisting}

Sin \texttt{fork}/\texttt{exec} separados, habría que pasarle a
\texttt{ls} (o a una hipotética \texttt{spawn}) un parámetro para
indicar la redirección. UNIX lo hace genérico: el shell modifica los
descriptores y \texttt{ls} simplemente escribe a \texttt{stdout}.

\paragraph{2) Pipes entre comandos.}

\begin{lstlisting}[style=cstyle]
/* `ls | wc -l` */
int p[2]; pipe(p);
if (fork() == 0) {        /* hijo izquierdo: ls */
    dup2(p[1], 1);
    close(p[0]); close(p[1]);
    execlp("ls", "ls", NULL);
}
if (fork() == 0) {        /* hijo derecho: wc */
    dup2(p[0], 0);
    close(p[0]); close(p[1]);
    execlp("wc", "wc", "-l", NULL);
}
close(p[0]); close(p[1]);
/* el shell espera ambos */
\end{lstlisting}

\paragraph{3) Cambio de ambiente.}
El hijo puede modificar variables de entorno con \texttt{setenv} o
construirse un \texttt{envp} propio antes de \texttt{exec}, sin
afectar al padre.

\paragraph{4) Cambio de identidad y privilegios.}
Programas SUID-root típicamente hacen \texttt{setuid()} para bajar
sus privilegios antes de \texttt{exec}, ejecutando ya como el usuario
real, no como root.

\paragraph{5) Cambio de directorio de trabajo.}
\texttt{chdir(\ldots)} en el hijo antes de \texttt{exec} para que
el nuevo programa arranque en otro directorio.

\paragraph{6) Configuración de signal handlers / mask.}
Resetear handlers a default antes del \texttt{exec}.

\paragraph{7) Limites de recursos.}
\texttt{setrlimit} (memoria máxima, tiempo de CPU, etc.) antes del
\texttt{exec}.

\paragraph{8) Uso de \texttt{fork} sin \texttt{exec}.}
También permite el patrón inverso: crear procesos hijos que ejecutan
el \emph{mismo programa} (paralelismo, workers, servidores con
modelo \emph{forking}). El servidor \texttt{Apache} pre-fork es el
ejemplo clásico.

\subsection*{Comparación con \texttt{spawn()}}

Windows tiene \texttt{CreateProcess()}, un único syscall con
\emph{muchísimos} argumentos para configurar todo lo anterior. UNIX
prefiere la \textbf{composición de syscalls pequeñas y ortogonales}
(redirección con \texttt{dup2}, cambio de privilegios con
\texttt{setuid}, etc.), y la flexibilidad de poder usar
\emph{cualquier} syscall durante la ventana del hijo.

\begin{respuesta}
La separación en dos pasos permite, en el período entre
\texttt{fork} y \texttt{exec}, configurar libremente el ambiente del
hijo: redirigir descriptores, conectarse a pipes, cambiar
variables de entorno, bajar privilegios, cambiar el directorio de
trabajo, configurar signal handlers, etc. Es la base de cómo el
shell implementa redirecciones (\texttt{>}, \texttt{<}), pipes
(\texttt{|}), substitución de comandos, jobs en background y demás
features. Si \texttt{fork+exec} fueran un solo syscall, todo eso
debería pasarse como parámetros.
\end{respuesta}

% =====================================================================
\section{Ejercicio 4 --- Características de xv6}
% =====================================================================

\textbf{Enunciado.} Tildar las opciones correctas sobre xv6 (código
original).

\subsection*{4.a --- ``Soporta spinlock''}

\begin{respuesta}
\textbf{V}. xv6 implementa spinlocks en \texttt{spinlock.c}: la
estructura \texttt{struct spinlock} con las funciones
\texttt{acquire()} y \texttt{release()}. Usa la instrucción atómica
\texttt{xchg} en x86 (\texttt{amoswap.w.aq} en RISC-V) para
implementar el test-and-set. Es el mecanismo de exclusión mutua
fundamental del kernel.
\end{respuesta}

\subsection*{4.b --- ``Soporta monitores''}

\begin{respuesta}
\textbf{F}. xv6 \emph{no} ofrece la abstracción de \emph{monitor}
(que en Java/C\# combina un lock con variables de condición en una
única construcción del lenguaje). Lo que xv6 sí ofrece son las
piezas separadas: \texttt{spinlock} para exclusión mutua y
\texttt{sleep}/\texttt{wakeup} para condition-variables-like
(item 4.c), pero el programador del kernel las combina manualmente.
\end{respuesta}

\subsection*{4.c --- ``El gestor de procesos tiene funciones comparables a variables de condición''}

\begin{respuesta}
\textbf{V}. Las funciones \texttt{sleep(chan, lock)} y
\texttt{wakeup(chan)} de xv6 implementan \emph{exactamente} la
semántica de variables de condición:

\begin{itemize}[leftmargin=*,nosep]
    \item \texttt{sleep(chan, lock)}: libera \texttt{lock}, pone al
        proceso a dormir en \texttt{chan} (estado \texttt{SLEEPING}).
        Al despertar, vuelve a tomar el lock. Equivalente a
        \texttt{pthread\_cond\_wait}.
    \item \texttt{wakeup(chan)}: despierta a \emph{todos} los
        procesos dormidos en \texttt{chan} (pasa a
        \texttt{RUNNABLE}). Equivalente a
        \texttt{pthread\_cond\_broadcast}.
\end{itemize}

xv6 no llama a estas operaciones ``variables de condición''
explícitamente, pero la analogía es directa.
\end{respuesta}

\subsection*{4.d --- ``No ofrece ningún mecanismo de comunicación y sincronización entre procesos''}

\begin{respuesta}
\textbf{F}. xv6 \emph{sí} ofrece mecanismos:

\begin{itemize}[leftmargin=*,nosep]
    \item \textbf{Pipes}: \texttt{pipe(p)} crea pipes entre procesos
        emparentados.
    \item \textbf{Archivos}: pueden usarse para comunicar procesos.
    \item \textbf{Signals (kill)}: aunque limitadas, xv6 implementa
        \texttt{kill(pid)}.
    \item \textbf{\texttt{wait}/\texttt{exit}}: sincronización
        padre-hijo a través del estado de salida.
    \item \textbf{Locks compartidos en filesystem}: \texttt{ilock} para
        inodos compartidos.
\end{itemize}
\end{respuesta}

\subsection*{4.e --- ``Tiene mecanismo de detección de deadlock''}

\begin{respuesta}
\textbf{F}. xv6 \emph{no} detecta deadlocks. La política es
``ostrich algorithm'': si los locks del kernel se toman siempre en
el mismo orden global (convención del código), el deadlock no
ocurrirá; si el programador comete un error de ordenamiento, xv6
simplemente se cuelga.

No hay grafo de espera, no hay timeout en \texttt{acquire}, no hay
recuperación. La prevención se hace por disciplina, no por detección
automática.
\end{respuesta}

\subsection*{Resumen Ej.\ 4}

\begin{center}
\small
\begin{tabular}{@{}lll@{}}
\toprule
Ítem & Respuesta & Marcar como verdadera \\
\midrule
a    & V         & \checkmark \\
b    & F         &           \\
c    & V         & \checkmark \\
d    & F         &           \\
e    & F         &           \\
\bottomrule
\end{tabular}
\end{center}

% =====================================================================
\section{Ejercicio 5 --- Evitar deadlock en par de locks invertidos}
% =====================================================================

\textbf{Enunciado.} Proponer una solución al código dado para evitar
deadlock. Describir ventajas/desventajas.

\begin{lstlisting}[style=cstyle]
// module A                     // module B
acquire(&a.lock);                acquire(&b.lock);
// update a.data                 // update b.value
acquire(&b.lock);                acquire(&a.lock);
// update b.value                // update a.data
release(&b.lock);                release(&a.lock);
release(&a.lock);                release(&b.lock);
\end{lstlisting}

\subsection*{Análisis: por qué hay deadlock}

Los módulos $A$ y $B$ adquieren los mismos dos locks en
\textbf{orden inverso}: $A$ toma $a$ y luego $b$; $B$ toma $b$ y
luego $a$. Una traza con ambos ejecutándose en paralelo:

\begin{enumerate}[leftmargin=*,nosep]
    \item $A$ ejecuta \texttt{acquire(\&a.lock)}: obtiene $a$.
    \item $B$ ejecuta \texttt{acquire(\&b.lock)}: obtiene $b$.
    \item $A$ ejecuta \texttt{acquire(\&b.lock)}: bloquea (B tiene $b$).
    \item $B$ ejecuta \texttt{acquire(\&a.lock)}: bloquea (A tiene $a$).
\end{enumerate}

Ambos bloqueados esperando un lock que el otro retiene $\to$
\textbf{deadlock} (las 4 condiciones de Coffman se cumplen).

\subsection*{Solución 1 (recomendada): orden global de adquisición}

Imponer un \textbf{orden total} sobre los locks que ambos módulos
respeten. Por ejemplo, ``siempre tomar primero el lock con menor
dirección de memoria'', o ``siempre primero \texttt{a.lock} y luego
\texttt{b.lock}''.

\begin{lstlisting}[style=cstyle]
/* Versión corregida: ambos módulos toman a.lock ANTES que b.lock */

// module A                     // module B
acquire(&a.lock);                acquire(&a.lock);  // CAMBIA: primero a
acquire(&b.lock);                acquire(&b.lock);  // luego b

// update a.data                 // update a.data
// update b.value                // update b.value

release(&b.lock);                release(&b.lock);
release(&a.lock);                release(&a.lock);
\end{lstlisting}

\paragraph{Por qué funciona.} Con orden global, nunca pueden formar
una espera circular: el primero que tome \texttt{a.lock} es el que
después podrá pedir \texttt{b.lock}, y el otro espera al primero.
Como mucho hay una cadena lineal, jamás un ciclo.

\paragraph{Ventajas.}
\begin{itemize}[leftmargin=*,nosep]
    \item \textbf{Simple} de implementar y razonar.
    \item \textbf{Sin penalidad de performance}: los locks se toman
        normalmente, sólo cambia el orden.
    \item \textbf{Granularidad fina}: dos hilos que toquen partes
        disjuntas del estado pueden ejecutar concurrentemente
        (sólo se serializan cuando colisionan en los mismos locks).
\end{itemize}

\paragraph{Desventajas.}
\begin{itemize}[leftmargin=*,nosep]
    \item Requiere disciplina del programador: si en algún lugar
        del código alguien viola el orden, vuelve el deadlock.
    \item En sistemas grandes con muchos locks, mantener el orden
        global se complica.
\end{itemize}

\subsection*{Solución 2: un único lock global}

Sustituir los dos locks por uno solo, que proteja todo el estado.

\paragraph{Ventajas.}
\begin{itemize}[leftmargin=*,nosep]
    \item \textbf{Imposible} formar deadlock con un único lock.
    \item Trivial de implementar.
\end{itemize}

\paragraph{Desventajas.}
\begin{itemize}[leftmargin=*,nosep]
    \item \textbf{Pérdida de concurrencia}: dos hilos que tocaran
        partes disjuntas del estado ahora se serializan
        innecesariamente.
    \item En sistemas con alta contención, este lock se vuelve
        cuello de botella.
\end{itemize}

\subsection*{Solución 3: \texttt{trylock} con \emph{back-off}}

Usar \texttt{trylock} para el segundo lock; si falla, soltar el
primero y reintentar.

\begin{lstlisting}[style=cstyle]
while (1) {
    acquire(&a.lock);
    if (trylock(&b.lock) == 0) {  /* tomado */
        break;
    }
    release(&a.lock);
    /* opcional: sleep o spin un rato antes de reintentar */
}
/* aqui A tiene a y b */
\end{lstlisting}

\paragraph{Ventajas.}
\begin{itemize}[leftmargin=*,nosep]
    \item No requiere coordinación global de orden.
    \item Funciona en escenarios donde no se sabe a priori el
        conjunto de locks.
\end{itemize}

\paragraph{Desventajas.}
\begin{itemize}[leftmargin=*,nosep]
    \item Más complejo de implementar.
    \item Puede generar \textbf{livelock}: dos hilos que se ceden
        mutuamente y nunca progresan. Requiere \emph{back-off
        aleatorio} para mitigarlo.
    \item Lecturas/escrituras parciales en estructura compartida
        si no se cuida con \texttt{release} oportuno.
\end{itemize}

\subsection*{Comparación}

\begin{center}
\small
\begin{tabular}{@{}lp{3.5cm}p{3.5cm}p{3.5cm}@{}}
\toprule
& \textbf{Orden global} & \textbf{Un solo lock} & \textbf{Trylock + back-off} \\
\midrule
Complejidad        & Baja                       & Mínima                 & Alta \\
Concurrencia       & Buena (locks separados)   & Mala (serialización)   & Buena (intentos no bloquean) \\
Riesgo de deadlock & 0 si se respeta el orden  & 0                      & 0, pero hay livelock \\
Aplicabilidad      & Cuando los locks se conocen estáticamente & Siempre & Cuando los locks son dinámicos \\
\bottomrule
\end{tabular}
\end{center}

\begin{respuesta}
La solución estándar es \textbf{imponer un orden global de
adquisición}: ambos módulos toman primero \texttt{a.lock} y luego
\texttt{b.lock} (o el orden inverso, pero consistente). Esto elimina
la espera circular sin sacrificar concurrencia. La desventaja es que
exige disciplina del programador para mantener el orden a lo largo
de toda la base de código.
\end{respuesta}

% =====================================================================
\section{Ejercicio 6 --- Traducción de direcciones lógicas a físicas}
% =====================================================================

\textbf{Enunciado.} Dado el mapa de memoria mostrado, determinar las
direcciones físicas correspondientes a:
\textbf{a)} \texttt{0x00000320}
\textbf{b)} \texttt{0x00002FFA}
\textbf{c)} \texttt{0x00400200}.
También determinar el rango de direcciones lógicas que mapean al
\textbf{frame 3}.

\subsection*{Análisis: estructura del mapa}

El diagrama muestra paginación de \textbf{dos niveles} estilo x86 32
bits (4\,KB páginas):

\begin{itemize}[leftmargin=*,nosep]
    \item Page directory con entradas $0, 1, 2, \ldots$
    \item Page table 1 con dos entradas: \texttt{0x2000} $\to$
        frame 1 y \texttt{0xA000} $\to$ frame 3.
    \item Page table 2 con una entrada: \texttt{0x5000} $\to$
        frame 2.
\end{itemize}

Las flechas del diagrama indican:

\begin{itemize}[leftmargin=*,nosep]
    \item \texttt{pgdir[0]} $\to$ Page Table 1.
    \item \texttt{pgdir[1]} $\to$ Page Table 2.
    \item PT1 [0] = \texttt{0x2000} \texttt{| flags} $\to$ frame 1
        en físico \texttt{0x02000000}.
    \item PT1 [2] = \texttt{0xA000} \texttt{| flags} $\to$ frame 3
        en físico \texttt{0x0A000000}.
    \item PT2 [0] = \texttt{0x5000} \texttt{| flags} $\to$ frame 2
        en físico \texttt{0x05000000}.
\end{itemize}

\textbf{Codificación de direcciones} (32 bits, páginas de 4\,KB):

\begin{itemize}[leftmargin=*,nosep]
    \item bits 31--22: \texttt{pgdir} index (10 bits).
    \item bits 21--12: \texttt{pgtable} index (10 bits).
    \item bits 11--0: \texttt{offset} (12 bits).
\end{itemize}

\textbf{Interpretación del PTE}: el valor del PTE (e.g.,
\texttt{0x2000}) representa el \textbf{frame number} (PPN); la
dirección física del frame es PPN $\times$ 4\,KB. Así:
\texttt{0x2000} $\to$ físico \texttt{0x2000} $\times$
\texttt{0x1000} $=$ \texttt{0x02000000}.

\subsection*{6.a --- Dirección lógica \texttt{0x00000320}}

Descomposición:

\quad \texttt{0x00000320} $=$ \texttt{0000\_0000\_00}\ %
\texttt{00\_0000\_0000}\ \texttt{0011\_0010\_0000}

\begin{itemize}[leftmargin=*,nosep]
    \item \texttt{pgdir} = $0$
    \item \texttt{pgtable} = $0$
    \item \texttt{offset} = \texttt{0x320}
\end{itemize}

Camino: \texttt{pgdir[0]} $\to$ PT1; \texttt{PT1[0]} = \texttt{0x2000}
$\to$ frame en físico \texttt{0x02000000}.

\quad \textbf{Físico = $\texttt{0x02000000} + \texttt{0x320} =
\texttt{0x02000320}$.}

\subsection*{6.b --- Dirección lógica \texttt{0x00002FFA}}

Descomposición:

\quad \texttt{0x00002FFA} $=$ \texttt{0000\_0000\_00}\ %
\texttt{00\_0000\_0010}\ \texttt{1111\_1111\_1010}

\begin{itemize}[leftmargin=*,nosep]
    \item \texttt{pgdir} = $0$
    \item \texttt{pgtable} = $2$
    \item \texttt{offset} = \texttt{0xFFA}
\end{itemize}

Camino: \texttt{pgdir[0]} $\to$ PT1; \texttt{PT1[2]} = \texttt{0xA000}
$\to$ frame en físico \texttt{0x0A000000}.

\quad \textbf{Físico = $\texttt{0x0A000000} + \texttt{0xFFA} =
\texttt{0x0A000FFA}$.}

\subsection*{6.c --- Dirección lógica \texttt{0x00400200}}

Descomposición:

\quad \texttt{0x00400200} $=$ \texttt{0000\_0000\_01}\ %
\texttt{00\_0000\_0000}\ \texttt{0010\_0000\_0000}

\begin{itemize}[leftmargin=*,nosep]
    \item \texttt{pgdir} = $1$
    \item \texttt{pgtable} = $0$
    \item \texttt{offset} = \texttt{0x200}
\end{itemize}

Camino: \texttt{pgdir[1]} $\to$ PT2; \texttt{PT2[0]} = \texttt{0x5000}
$\to$ frame en físico \texttt{0x05000000}.

\quad \textbf{Físico = $\texttt{0x05000000} + \texttt{0x200} =
\texttt{0x05000200}$.}

\subsection*{6.d --- Rango lógico que mapea al frame 3}

\begin{importante}
\textbf{El rango original \texttt{[0x00400000, 0x0040BFFF]} estaba
incorrecto.} El razonamiento dado (``\texttt{pgdir[1]} +
\texttt{pgtable[0]}'') apuntaba al frame 2, no al frame 3.
\end{importante}

\textbf{Razonamiento correcto.} El frame 3 se alcanza a través de
\texttt{pgdir[0]} (que apunta a PT1) y \texttt{pgtable[2]} (entrada
\texttt{0xA000} de PT1). Esto es \emph{coherente} con la respuesta de
6.b: \texttt{0x00002FFA} mapeó a frame 3.

Por lo tanto, las direcciones lógicas que caen en frame 3 son
aquellas con:

\begin{itemize}[leftmargin=*,nosep]
    \item bits 31--22 (\texttt{pgdir}) $= 0$
    \item bits 21--12 (\texttt{pgtable}) $= 2$
    \item bits 11--0 (\texttt{offset}) $\in [0, \texttt{0xFFF}]$
\end{itemize}

Composición:

\quad mínimo: \texttt{0000\_0000\_00\_00\_0000\_0010\_0000\_0000\_0000}
$=$ \texttt{0x00002000}

\quad máximo: \texttt{0000\_0000\_00\_00\_0000\_0010\_1111\_1111\_1111}
$=$ \texttt{0x00002FFF}

\begin{respuesta}
\textbf{Direcciones físicas:}

\begin{center}
\small
\begin{tabular}{@{}llll@{}}
\toprule
Lógica           & \texttt{pgdir} & \texttt{pgtable[idx]} & Física \\
\midrule
\texttt{0x00000320} & 0 & PT1[0] = 0x2000 & \textbf{\texttt{0x02000320}} \\
\texttt{0x00002FFA} & 0 & PT1[2] = 0xA000 & \textbf{\texttt{0x0A000FFA}} \\
\texttt{0x00400200} & 1 & PT2[0] = 0x5000 & \textbf{\texttt{0x05000200}} \\
\bottomrule
\end{tabular}
\end{center}

\textbf{Rango lógico del frame 3}: \boxed{\texttt{[0x00002000,~0x00002FFF]}}
(4\,KB consecutivos, accedidos vía \texttt{pgdir[0] + pgtable[2]}).
\end{respuesta}

% =====================================================================
\section{Ejercicio 7 --- Diagrama del mapa de memoria de un proceso xv6}
% =====================================================================

\textbf{Enunciado.} Diagrame el mapa (directorio + tablas de páginas)
para un proceso xv6 que usa 15\,KB de código, 3\,KB de datos
(y el stack correspondiente).

\subsection*{Análisis: layout que produce \texttt{exec} en xv6}

xv6 (versión del libro) coloca en el espacio de direcciones del
proceso, comenzando en \texttt{0}:

\begin{enumerate}[leftmargin=*,nosep]
    \item Las páginas de \textbf{código (\texttt{text})}.
    \item Las páginas de \textbf{datos} (\texttt{data} + \texttt{bss}).
    \item Una \textbf{página de guarda} con \texttt{V = 0} (para
        detectar \emph{stack overflow}).
    \item Una \textbf{página de stack} (el código de \texttt{exec}
        en xv6 asigna 1 página de stack utilizable, en
        \texttt{exec.c}: \emph{``Allocate two pages \ldots\ make the
        first inaccessible \ldots\ use the second as the user
        stack''}).
\end{enumerate}

\subsection*{Cuentas previas}

\begin{itemize}[leftmargin=*,nosep]
    \item \textbf{Código (15\,KB)}: $\lceil 15 / 4 \rceil = 4$
        páginas (la última con 1\,KB de \emph{padding}).
        $\to$ vpns \textbf{0, 1, 2, 3}.
    \item \textbf{Datos (3\,KB)}: $\lceil 3 / 4 \rceil = 1$
        página (con 1\,KB sin usar). $\to$ vpn \textbf{4}.
    \item \textbf{Guarda}: $1$ página con \texttt{V = 0}. $\to$ vpn
        \textbf{5}.
    \item \textbf{Stack}: $1$ página. $\to$ vpn \textbf{6}.
\end{itemize}

Total: \textbf{7 vpns}, de las cuales \textbf{6 están mapeadas}
(la guarda es inválida). Todas caen en los primeros 28\,KB del
espacio de direcciones; todas pertenecen al primer entrada del page
directory (\texttt{pgdir[0]}, que cubre los primeros 4\,MB).

\subsection*{Diagrama del directorio + tabla de páginas}

\begin{center}
\begin{tikzpicture}[
    every node/.style={font=\ttfamily\small},
    pde/.style={draw,fill=c1!25,minimum width=3.2cm,minimum height=0.55cm,align=center},
    pte/.style={draw,fill=c3!20,minimum width=4cm,minimum height=0.5cm,align=center},
    inv/.style={draw,fill=cred!15,minimum width=4cm,minimum height=0.5cm,align=center,font=\ttfamily\scriptsize},
    page/.style={draw,minimum width=2.6cm,minimum height=0.5cm,align=center}
]

% Page directory
\node[font=\bfseries\small] at (-3.7,6) {Page Directory};
\node[pde] (pd0) at (-3.7,5.3) {[0] $\to$ \texttt{PT0}};
\node[inv,minimum width=3.2cm] (pdrest) at (-3.7,4.75) {[1]..[511] $V=0$};
\node[inv,minimum width=3.2cm,fill=c4!15] (pdkern) at (-3.7,4.2) {[512]..[1023] kernel};

% Page Table 0
\node[font=\bfseries\small] at (2,6) {Page Table 0 (\texttt{PT0})};
\node[pte] (p0) at (2,5.3) {[0] $\to$ F$_0$ (code, R+X)};
\node[pte,below=0pt of p0] (p1) {[1] $\to$ F$_1$ (code, R+X)};
\node[pte,below=0pt of p1] (p2) {[2] $\to$ F$_2$ (code, R+X)};
\node[pte,below=0pt of p2] (p3) {[3] $\to$ F$_3$ (code, R+X)};
\node[pte,below=0pt of p3,fill=c2!25] (p4) {[4] $\to$ F$_4$ (data, R+W)};
\node[inv,below=0pt of p4] (p5) {[5] (guarda, $V=0$)};
\node[pte,below=0pt of p5,fill=c4!25] (p6) {[6] $\to$ F$_5$ (stack, R+W)};
\node[inv,below=0pt of p6] (prest) {[7]..[1023] $V=0$};

% Frames físicos
\node[font=\bfseries\small] at (8.5,6) {Frames físicos};
\node[page,fill=c3!20] (f0) at (8.5,5.3) {F$_0$ (4\,KB)};
\node[page,fill=c3!20,below=0pt of f0] (f1) {F$_1$ (4\,KB)};
\node[page,fill=c3!20,below=0pt of f1] (f2) {F$_2$ (4\,KB)};
\node[page,fill=c3!20,below=0pt of f2] (f3) {F$_3$ (4\,KB)};
\node[page,fill=c2!25,below=0pt of f3] (f4) {F$_4$ (4\,KB)};
\node[page,fill=c4!25,below=0.4cm of f4] (f5) {F$_5$ (4\,KB)};

\draw[->,thick] (pd0.east) -- (p0.west);

\draw[->,thick] (p0.east) -- (f0.west);
\draw[->,thick] (p1.east) -- (f1.west);
\draw[->,thick] (p2.east) -- (f2.west);
\draw[->,thick] (p3.east) -- (f3.west);
\draw[->,thick] (p4.east) -- (f4.west);
\draw[->,thick] (p6.east) -- (f5.west);

\end{tikzpicture}
\end{center}

\subsection*{Tabla detallada de entradas válidas}

\begin{center}
\small
\begin{tabular}{@{}cclccp{3.5cm}@{}}
\toprule
vpn & rango virtual              & PT0[idx]   & V & permisos & Área \\
\midrule
0   & \texttt{[0x0000, 0x1000)}  & F$_0$      & 1 & R+X      & código (pág.\ 1) \\
1   & \texttt{[0x1000, 0x2000)}  & F$_1$      & 1 & R+X      & código (pág.\ 2) \\
2   & \texttt{[0x2000, 0x3000)}  & F$_2$      & 1 & R+X      & código (pág.\ 3) \\
3   & \texttt{[0x3000, 0x4000)}  & F$_3$      & 1 & R+X      & código (pág.\ 4, 1\,KB padding) \\
4   & \texttt{[0x4000, 0x5000)}  & F$_4$      & 1 & R+W      & datos (3\,KB + 1\,KB pad) \\
5   & \texttt{[0x5000, 0x6000)}  & ---        & 0 & ---      & \textbf{guarda (stack overflow)} \\
6   & \texttt{[0x6000, 0x7000)}  & F$_5$      & 1 & R+W      & stack \\
\bottomrule
\end{tabular}
\end{center}

El resto del \texttt{Page Table 0} (entradas 7 a 1023), así como el
resto del \texttt{Page Directory} (entradas 1 a 511 para usuario),
quedan con \texttt{V = 0}. Las entradas 512 a 1023 del PD están
mapeadas al kernel (\texttt{KERNBASE} $=$ 2\,GB en xv6 x86).

\begin{respuesta}
El proceso usa \textbf{6 frames físicos} (4 de código, 1 de datos,
1 de stack) más \textbf{1 página de guarda} (sin frame asociado).
Todo el layout cabe en los primeros 28\,KB del espacio virtual:
una sola entrada del \texttt{Page Directory} (\texttt{PD[0]}) que
apunta a un único \texttt{Page Table} (\texttt{PT0}) con 7 entradas
relevantes (de las cuales una es inválida).
\end{respuesta}

% =====================================================================
\section{Ejercicio 8 --- Segmentación y paginado}
% =====================================================================

\textbf{Enunciado.} Marcar las respuestas verdaderas.

\subsection*{8.a --- ``Los segmentos representan porciones lógicas de un proceso''}

\begin{respuesta}
\textbf{V}. Es justamente la idea fundamental de la segmentación:
dividir el espacio de direcciones de un proceso en porciones
\emph{con significado lógico} (código, datos, stack, heap, librería
compartida\ldots), cada una con su propio tamaño, base y permisos.
\end{respuesta}

\subsection*{8.b --- ``Un segmento o una página sólo puede contener instrucciones o datos, nunca ambos''}

\begin{respuesta}
\textbf{F}. Una \textbf{página} no distingue ``instrucciones de
datos'': es una unidad puramente física (4\,KB de memoria). Una sola
página puede contener instrucciones y datos mezclados (de hecho, el
último byte de un binario regularmente combina código y datos en
la misma página por alineamiento).

En \textbf{segmentación pura}, sí hay segmentos especializados
(segmento de código vs.\ segmento de datos), pero incluso ahí el
segmento de código puede contener constantes (\texttt{.rodata}). La
afirmación absoluta es falsa.
\end{respuesta}

\subsection*{8.c --- ``El paginado es comúnmente usado para implementar memoria virtual''}

\begin{respuesta}
\textbf{V}. La paginación es \emph{la} técnica estándar para
memoria virtual en SO modernos. Sus ventajas: páginas de tamaño
fijo (swap simple), sin fragmentación externa, soporte natural para
demand paging y copy-on-write.
\end{respuesta}

\subsection*{8.d --- ``En x86 las páginas pueden ser de un nivel o de dos niveles''}

\begin{respuesta}
\textbf{V}. x86 32-bit admite dos modos de paginación:

\begin{itemize}[leftmargin=*,nosep]
    \item \textbf{2 niveles}: PD + PT, páginas de \textbf{4\,KB}. La
        configuración estándar.
    \item \textbf{1 nivel}: sólo el PD con páginas de \textbf{4\,MB}
        (extensión PSE, \emph{Page Size Extension}). Cada entrada del
        PD apunta directamente a una página de 4\,MB. Útil para
        mapear regiones grandes del kernel.
\end{itemize}

x86 64-bit usa 4 ó 5 niveles, pero en el contexto de x86 32-bit
(que es lo que cubre la cátedra), la afirmación es V.
\end{respuesta}

\subsection*{8.e --- ``En xv6 (original) al ocurrir un page fault, el proceso finaliza (\texttt{exit})''}

\begin{respuesta}
\textbf{V}. xv6 \emph{no implementa} demand paging, COW ni
crecimiento dinámico del stack. Cualquier page fault en modo usuario
es tratado como un acceso ilegal: \texttt{trap.c} marca
\texttt{p->killed = 1}, lo que provoca que el proceso termine en su
próximo \texttt{trapret}.

En SO más sofisticados (Linux), un page fault es un \emph{mecanismo
útil} para implementar varias optimizaciones; sólo se mata el
proceso si la dirección no pertenece a ningún VMA. Pero en xv6
original, \textbf{todo} page fault de usuario es fatal.
\end{respuesta}

\subsection*{Resumen Ej.\ 8}

\begin{center}
\small
\begin{tabular}{@{}lll@{}}
\toprule
Ítem & Respuesta & Marcar como verdadera \\
\midrule
a    & V         & \checkmark \\
b    & F         &           \\
c    & V         & \checkmark \\
d    & V         & \checkmark \\
e    & V         & \checkmark \\
\bottomrule
\end{tabular}
\end{center}

% =====================================================================
\section{Ejercicio 9 --- Una sola tabla de páginas en x86 para 4 GB}
% =====================================================================

\textbf{Enunciado.} Si se usa una única tabla de páginas en x86
(32 bits). Para describir 4\,GB se requiere que:

\begin{itemize}[leftmargin=*,nosep]
    \item[(a)] sus entradas describan páginas de \textbf{4\,KB}.
    \item[(b)] sus entradas describan páginas de \textbf{4\,MB}.
\end{itemize}

\subsection*{Análisis}

``Una única tabla de páginas'' significa paginación de \emph{un solo
nivel} (sin Page Directory + Page Table; sólo una tabla plana).

Para cubrir los $2^{32}$ bytes (4\,GB) del espacio de direcciones,
las entradas de la tabla y el offset deben sumar 32 bits. Si cada
entrada cubre una página de tamaño $P$, hace falta una tabla con
$2^{32} / P$ entradas.

\textbf{Caso (a) páginas de 4\,KB:} la tabla tendría
$2^{32} / 2^{12} = 2^{20} = 1\,048\,576$ entradas. Cada entrada
ocupa $4$ bytes $\to$ \textbf{$4$\,MB} de tabla por proceso. Además,
en x86 esto no se hace ``en una sola tabla''; el hardware del
procesador requiere PD + PT (2 niveles) para páginas de 4\,KB.

\textbf{Caso (b) páginas de 4\,MB:} la tabla tendría
$2^{32} / 2^{22} = 2^{10} = 1024$ entradas. Cada entrada ocupa
$4$ bytes $\to$ \textbf{$4$\,KB} de tabla por proceso (cabe en una
sola página). \textbf{x86 lo soporta de forma nativa} mediante el
bit PS (Page Size) en las entradas del Page Directory: si está
encendido, la entrada apunta directamente a una página de 4\,MB en
lugar de a otro nivel.

\begin{respuesta}
La respuesta correcta es \textbf{(b) páginas de 4\,MB}. Una única
tabla de 1024 entradas (4\,KB de tamaño) cubre los 4\,GB con
páginas de 4\,MB cada una. Páginas de 4\,KB exigirían 1\,M de
entradas (4\,MB de tabla), y x86 no soporta esa configuración con
una sola tabla; requiere paginación de dos niveles.
\end{respuesta}

% =====================================================================
\section{Ejercicio 10 --- Concepto de memoria virtual}
% =====================================================================

\textbf{Enunciado.} Describir el concepto de memoria virtual.

\subsection*{Definición}

La \textbf{memoria virtual} es una abstracción provista
conjuntamente por el hardware (MMU) y el SO que desacopla la
\emph{dirección lógica} (la que un proceso emite cuando lee o
escribe memoria) de la \emph{dirección física} (la que en realidad
selecciona una celda de la RAM).

Esto permite que cada proceso tenga su \textbf{propio espacio de
direcciones virtuales} independiente, que puede ser \emph{más
grande} que la RAM física disponible y que el SO administre
dinámicamente.

\subsection*{Mecanismo básico}

\begin{enumerate}[leftmargin=*,nosep]
    \item Cada proceso tiene su propia \textbf{tabla de páginas}
        que el SO mantiene.
    \item Cuando el proceso emite una dirección virtual, la
        \textbf{MMU} la consulta en la tabla y la traduce a la
        dirección física correspondiente.
    \item Si la página no está en RAM (\texttt{V = 0}), la MMU
        genera un \emph{page fault}. El SO atiende el fault, decide
        qué hacer (cargar la página de disco, asignar una nueva,
        copiar, matar al proceso, etc.) y reanuda la instrucción.
\end{enumerate}

\subsection*{Funcionalidades que habilita}

\paragraph{1) Procesos más grandes que la RAM.}
\textbf{Swapping}: páginas no usadas recientemente se mueven a una
zona de \emph{swap} en disco. Cuando se vuelven a referenciar, el
SO las trae de vuelta (vía page fault). Permite ejecutar procesos
que en conjunto exceden la RAM, al costo de latencia de disco.

\paragraph{2) Demand paging.}
Las páginas se cargan en RAM \textbf{sólo cuando se accede a ellas
por primera vez}. Esto acelera el arranque de programas grandes:
no hay que leer el ejecutable entero al disco al hacer \texttt{exec},
sólo las páginas que se ejecutan.

\paragraph{3) Aislamiento entre procesos.}
Como cada proceso ve solo su tabla de páginas, no puede acceder a la
memoria de otro. Es la base de la seguridad multi-usuario.

\paragraph{4) Copy-on-Write (COW).}
\texttt{fork} no copia la memoria del padre. Las páginas se
comparten físicamente con permiso de \texttt{lectura}; el primer
\texttt{write} dispara un page fault que el kernel resuelve copiando
sólo la página afectada. Hace \texttt{fork+exec} extremadamente
eficiente.

\paragraph{5) Memory-mapped files (\texttt{mmap}).}
Un archivo se ``mapea'' al espacio de direcciones del proceso. Las
páginas se cargan bajo demanda; lecturas y escrituras se traducen
automáticamente en operaciones sobre el archivo. Es la base de
cargar ejecutables eficientemente.

\paragraph{6) Memoria compartida entre procesos.}
Dos procesos pueden mapear las mismas páginas físicas a (posiblemente
distintas) direcciones virtuales. Es la base de \texttt{shm}, IPC
de alto rendimiento y librerías compartidas.

\paragraph{7) Páginas de guarda.}
El SO puede dejar páginas con \texttt{V = 0} entre regiones críticas
(e.g., entre el stack y el heap) para detectar overflows
inmediatamente.

\paragraph{8) Address Space Layout Randomization (ASLR).}
La MMU permite cargar el binario en direcciones aleatorias en cada
ejecución, dificultando ataques que dependen de direcciones
conocidas.

\subsection*{Costos}

\begin{itemize}[leftmargin=*,nosep]
    \item \textbf{Overhead de traducción}: cada acceso pasa por la
        MMU. Para mitigar, las CPUs tienen un \textbf{TLB} (cache
        de traducciones).
    \item \textbf{Page faults} cuestan miles de ciclos. Si se
        dispara mucho (\emph{thrashing}, working set $>$ RAM), el
        sistema se cae a pique.
    \item \textbf{Memoria de la tabla de páginas} (algunos MB por
        proceso en x86-64), aunque el árbol jerárquico mitiga
        cuando hay grandes huecos.
\end{itemize}

\begin{respuesta}
\textbf{Memoria virtual} es la abstracción que da a cada proceso un
espacio de direcciones lógicas independiente, posiblemente mayor
que la RAM física, gestionado por el SO con ayuda de la MMU. Sus
funcionalidades clave: aislamiento de procesos, swapping,
\emph{demand paging}, \emph{copy-on-write}, \emph{memory-mapped
files}, memoria compartida entre procesos, ASLR y páginas de guarda.

La idea unificadora: la tabla de páginas es una indirección
flexible entre lo que el proceso ``ve'' y lo que la RAM ``es'', y
el SO la usa para implementar todas las features anteriores con un
único mecanismo (resolver el page fault apropiadamente).
\end{respuesta}

% =====================================================================
\section*{Resumen del parcial}
\addcontentsline{toc}{section}{Resumen del parcial}
% =====================================================================

\begin{center}
\footnotesize
\begin{tabular}{@{}clp{9cm}@{}}
\toprule
\# & Tema & Respuesta sintética \\
\midrule
1  & Diseño de SO        & a:V, b:F, c:V, \textbf{d:V} (microkernel distribuible), e:F \\
2  & Gestión de procesos & a:V, b:F (zombie ya finalizó), c:F (exec no crea), \textbf{d:F}, \textbf{e:F} \\
3  & \texttt{fork} + \texttt{exec} & Ventana para configurar el ambiente del hijo: redirecciones, pipes, env, privilegios, dir.\ de trabajo, signal handlers \\
4  & xv6                 & a:V (spinlock), b:F (sin monitor), c:V (\texttt{sleep}/\texttt{wakeup}), d:F, e:F \\
5  & Solución a deadlock & Orden global de adquisición; alternativas: lock único o trylock+backoff \\
6  & Traducción de direcciones & a:\texttt{0x02000320}, b:\texttt{0x0A000FFA}, c:\texttt{0x05000200}; frame 3: \textbf{\texttt{[0x00002000, 0x00002FFF]}} \\
7  & Mapa xv6 (15K + 3K) & PD[0] $\to$ PT0; PT0[0..3]=code, [4]=data, [5]=guarda (V=0), [6]=stack \\
8  & Segmentación/paginado & a:V, b:F, c:V, d:V, e:V \\
9  & Tabla única para 4\,GB & b: páginas de 4\,MB (1024 entradas, 4\,KB de tabla) \\
10 & Memoria virtual     & Abstracción para aislamiento, swap, demand paging, COW, mmap, shared mem, ASLR, guard pages \\
\bottomrule
\end{tabular}
\end{center}

\begin{importante}
\textbf{Correcciones respecto de las respuestas originales:}

\begin{enumerate}[leftmargin=*,nosep]
    \item \textbf{Ej 1.d}: era \textbf{V} (microkernel distribuible),
        no ``no sé''.
    \item \textbf{Ej 2.d, 2.e}: ambas \textbf{F}; estaban
        incompletas.
    \item \textbf{Ej 4.b}: \textbf{F} (xv6 no tiene monitores como
        construcción explícita; tiene los componentes por separado).
    \item \textbf{Ej 6.d (rango del frame 3)}: corregido a
        \texttt{[0x00002000, 0x00002FFF]}. El rango original
        \texttt{[0x00400000, 0x0040BFFF]} era inconsistente con las
        respuestas (a), (b) y (c).
\end{enumerate}
\end{importante}

\end{document}
